%!TEX root = ../main.tex
% _Trim behavious: The action of adjusting control surfaces to maintain stable flight.
%_Reduced order aerodynamics analyses: Simplified methods for predicting aerodynamic forces and moments on airfoils, often using empirical data or simplified mathematical models.
% Parametric variations: Systematic changes in key parameters (e.g., angle of attack, Reynolds number) to study their effects on aerodynamic performance.
% Abstract to be finished later.

\begin{abstract}
\noindent A ground-effect vehicle (GEV) is an aircraft designed to operate in close proximity to the ground, capitalising on a phenomenon -- the so-called ground effect -- that results in an significant increase in lift and a  reduction in induced drag. Despite their potential efficiency advantages, GEVs present significant challenges in airfoil selection and trim behavior, particularly when transitioning between ground-effect and free-flight conditions.\\

\noindent This study investigates the aerodynamic characteristics of selected airfoil profiles in ground effect, with the objective of identifying configurations suitable for mechanically stable GEV operation. Building upon preliminary airfoil investigations conducted during a prior academic project, \textbf{five} representative airfoils are selected and evaluated using a combination of reduced-order aerodynamic analysis and computational fluid dynamics (CFD) simulations. Parametric variations in angle of attack, Reynolds number, and height-to-chord ratio are considered to assess lift, drag, and pitching moment behavior.\\
\end{abstract}

